\section{Symmetric groups}\label{sec:symmetric_groups}

\paragraph{Symmetric groups and permutations}

\begin{definition}\label{def:symmetric_group}\mcite[def. II.2.1]{Aluffi2009Algebra}
  We call the \hyperref[def:automorphism_group]{automorphism group} of a set \( A \) the \term[bg=симетрична група (\cite[376]{Обрешков1962ВисшаАлгебра}), ru=симметрическая группа (\cite[154]{Винберг2014КурсАлгебры})]{symmetric group}\fnote{The term is closely related to symmetric functions defined in \cref{def:symmetric_function}} on \( A \) and denote it by \( S_A \). The group \( S_A \) consists of bijective functions, which we call \term[bg=субституция (\cite[80]{ГеновМиховскиМоллов1991Алгебра}), ru=подстановка (\cite[154]{Винберг2014КурсАлгебры})); перестановка (\cite[sec. 4.2]{Тыртышников2007ЛинейнаяАлгебра})]{permutations}.

  For a \hyperref[def:cardinal]{cardinal number} \( \nu \), which by \cref{thm:ordinal_is_set_of_smaller_ordinals} is the set of all smaller ordinals, \( S_\nu \) acts on these smaller ordinals. This is similar to using cardinals for colorings discussed in \cref{rem:cardinal_colorings} or vertices on universal graphs discussed in \cref{rem:universal_graph_cardinals}. By convention, however, for a natural number \( n \), the group \( S_n \) to act on \( \set{ 1, \ldots, n } \) rather than \( \set{ 0, \ldots, n - 1 } \). We can denote the permutation \( \sigma \) in \( S_n \) as
  \begin{equation*}
    \sigma
    =
    \begin{pmatrix}
      1         & \cdots & n \\
      \sigma(1) & \cdots & \sigma(n)
    \end{pmatrix}.
  \end{equation*}
\end{definition}
\begin{comments}
  \item Some authors, for example \textcite[253]{MacLane1998CategoryTheory}, \textcite[31]{Jacobson1985BasicAlgebraI} and \textcite[121]{Knapp2016BasicAlgebra}, call \( S_n \) the \enquote{symmetric group on \( n \) letters}.

  \item The term \enquote{permutation} is also used in combinatorics, where it has a different meaning that only coincides with this usage in the cases discussed in \cref{thm:combinatorial_and_algebraic_permutation}. See \cref{rem:combinatorial_permutation_terminology} for a discussion of the term \enquote{permutation}.

  \item Our first complete example of a symmetric group will be given for \( S_3 \) in \cref{ex:s3}, since we need additional machinery in order to properly study it.

  \item \Cref{thm:symmetric_group_cardinality} implies that \( S_n \) has \( n! \) elements.
\end{comments}

\begin{theorem}[Cayley's theorem]\label{thm:cayleys_theorem}
  Every \hyperref[def:group]{group} \( G \) \hyperref[def:fol_embedding]{embeds isomorphically} into the corresponding \hyperref[def:symmetric_group]{symmetric group} \( S_G \).
\end{theorem}
\begin{comments}
  \item Compare this to \fullref{thm:cayleys_theorem_for_monoids}.
\end{comments}
\begin{proof}
  The embedding is given by the group action from \cref{thm:group_is_action}.
\end{proof}

\begin{example}\label{ex:def:symmetric_group}
  We list examples of \hyperref[def:symmetric_group]{symmetric groups and permutations}
  \begin{thmenum}
    \thmitem{ex:def:symmetric_group/automorphism} Every \hyperref[def:morphism_invertibility/automorphism]{automorphism} is a permutation, but not vice versa.

    For example, the additive group of integers \( \BbbZ \) has an automorphism given by \( x \mapsto -x \), and this is a permutation since the only requirement for being a permutation is being a bijective function.

    But an arbitrary permutation like \( x \mapsto x + 1 \) is not an automorphism since \( 0 \mapsto 1 \).

    \thmitem{ex:def:symmetric_group/diamond} The permutation
    \begin{equation}\label{eq:ex:def:symmetric_group/diamond}
      \begin{pmatrix}
        1 & 2 & 3 & 4 & 5 & 6 \\
        3 & 4 & 5 & 6 & 1 & 2
      \end{pmatrix}
    \end{equation}
    is displayed graphically in \cref{fig:ex:def:symmetric_group/diamond}.

    The permutation sends even numbers to even numbers and odd numbers to odd numbers. In fact, we can represent it as the composition
    \begin{equation*}
      \underbrace
        {
          \begin{pmatrix}
            \mathbf{1} & 2 & \mathbf{3} & 4 & \mathbf{5} & 6 \\
            \mathbf{3} & 2 & \mathbf{5} & 4 & \mathbf{1} & 6
          \end{pmatrix}
        }
        _
        {
          \cycle{ 1, 3, 5 }
        }
      \bincirc
      \underbrace
        {
          \begin{pmatrix}
            1 & \mathbf{2} & 3 & \mathbf{4} & 5 & \mathbf{6} \\
            1 & \mathbf{4} & 3 & \mathbf{6} & 5 & \mathbf{2}
          \end{pmatrix}.
        }
        _
        {
          \cycle{ 2, 4, 6 }
        }
    \end{equation*}

    The shortened notation in the underbraces is introduced in \cref{rem:cycle_notation} and, for these concrete permutations, discussed in \cref{ex:def:cyclic_permutation/diamond}.

    \begin{figure}
      \centering
      \includegraphics{output/ex__def__symmetric_group}
      \caption{A visualization of the permutation from \cref{ex:def:symmetric_group/diamond}}
      \label{fig:ex:def:symmetric_group/diamond}
    \end{figure}
  \end{thmenum}
\end{example}

\begin{lemma}\label{thm:sum_of_powers_in_composition}
  For any endomorphism \( f: A \to A \) in any \hyperref[def:category]{category}, for any integers \( n \) and \( m \) we have
  \begin{equation*}
    f^{n + m} = f^n \bincirc f^m.
  \end{equation*}
\end{lemma}
\begin{proof}
  Follows via simple induction on \( m \).
\end{proof}

\begin{proposition}\label{thm:symmetric_group_action}
  Given any \hyperref[def:symmetric_group]{permutation} \( \sigma \) from \( S_A \), the additive group of integers \( \BbbZ \) \hyperref[def:group_action]{acts} on \( A \) via
  \begin{equation*}
    \Phi_e(a) \coloneqq \sigma^e(a).
  \end{equation*}
\end{proposition}
\begin{proof}
  \Cref{thm:sum_of_powers_in_composition} implies that the condition \eqref{eq:def:group_action/family/compatibility} holds.
\end{proof}

\begin{remark}\label{rem:symmetric_group_actions}
  The symmetric group \( S_A \) acts on \( A \) via
  \begin{equation*}
    \Phi_\sigma(a) \coloneqq \sigma(a).
  \end{equation*}

  This should not be confused with the action \cref{thm:symmetric_group_action}, with respect to which we consider orbits in \cref{def:cyclic_permutation}.
\end{remark}

\begin{proposition}\label{thm:symmetric_group_universal_property}
  Every function \( f: A \to B \) extends to a \hyperref[def:group/morphism]{homomorphism} \( \widetilde{f}: S_A \to S_B \) between their \hyperref[def:symmetric_group]{symmetric groups}.

  Furthermore, if \( f \) is injective, so is \( \widetilde{f} \).
\end{proposition}
\begin{proof}
  Let
  \begin{equation*}
    \widetilde{f}(\sigma)(y) \coloneqq \begin{cases}
      f(\sigma(x)), &y = f(x), \\
      y,            &\T{otherwise.}
    \end{cases}
  \end{equation*}

  Then \( \widetilde{f} \) is a homomorphism because
  \begin{equation*}
    \widetilde{f}(\pi \bincirc \sigma)(y)
    =
    \begin{rcases}
      \begin{cases}
        f(\pi(\sigma(x))), &y = f(x), \\
        y,                 &\T{otherwise}
      \end{cases}
    \end{rcases}
    =
    [\widetilde{f}(\pi) \bincirc \widetilde{f}(\sigma)](y)
  \end{equation*}

  Furthermore, if \( f \) is injective and \( \widetilde{f}(\sigma) = \widetilde{f}(\pi) \), then \( \sigma = \pi \). Indeed, for every \( x \) in \( A \), we have \( f(\sigma(x)) = f(\pi(x)) \), and thus \( \sigma(x) = \pi(x) \).
\end{proof}

\begin{corollary}\label{thm:sn_embeds_into_sm}
  There are \( n! / (n - k)! \) \hyperref[def:group/morphism]{embeddings} of the \hyperref[def:symmetric_group]{symmetric group} \( S_k \) into \( S_n \) --- one for every injective function from \( \set{ 1, \ldots, k } \) to \( \set{ 1, \ldots, n } \).
\end{corollary}
\begin{comments}
  \item In particular, embeddings of \( S_2 \) correspond exactly to \hyperref[def:transposition]{transpositions}.
\end{comments}
\begin{proof}
  That the functions induce homomorphisms follows from \cref{thm:symmetric_group_universal_property}, while the number of embeddings can be counted via \cref{thm:combinatorial_variation_count/no_repetition}.
\end{proof}

\paragraph{Cyclic permutations}

\begin{definition}\label{def:cyclic_permutation}\mcite[def. II.4.1]{Aluffi2009Algebra}
  We say that the \hyperref[def:symmetric_group]{permutation} \( \sigma \) in \( S_n \) is \term{cyclic} or a \term[ru=цикл (\cite[sec. 4.3]{Тыртышников2007ЛинейнаяАлгебра})]{cycle} if the action \( \Phi_e \) from \cref{thm:symmetric_group_action} has \emph{at most} one \hyperref[def:group_action_orbit]{orbit} that is nontrivial (i.e. has more than one element).

  We call the cardinality of this orbit the \term{length} of the cycle. The identity can be viewed as a degenerate cycle of length \( 1 \) since it has only trivial orbits.
\end{definition}
\begin{comments}
  \item Cyclic permutations are defined only for the finite symmetric groups \( S_n \).
  \item We elaborate on the definition in \cref{thm:cyclic_permutation_characterization} and introduce a special cycle notation in \cref{rem:cycle_notation}.
\end{comments}

\begin{proposition}\label{thm:cyclic_permutation_characterization}
  A \hyperref[def:cyclic_permutation]{cyclic permutation} \( \sigma \) in \( S_n \) can be characterized as follows:
  \begin{thmenum}
    \thmitem{thm:cyclic_permutation_characterization/cycle} Pick one element from the nontrivial orbit and denote it by \( s_1 \) (if all orbits are trivial, pick any element). Recursively define
    \begin{equation*}
      s_{k+1} \coloneqq \sigma(s_k).
    \end{equation*}
    for all \( k < m \). Then
    \begin{equation}\label{eq:thm:def:cyclic_permutation/cycle}
      \sigma(q) = \begin{cases}
        s_1,       &q = s_m, \\
        s_{i + 1}, &q = s_i \T{for} i < m, \\
        q,         &\T{otherwise}
      \end{cases}
    \end{equation}

    \begin{figure}
      \centering
      \includegraphics{output/thm__cyclic_permutation_characterization}
      \caption{Visualization of a cycle with five elements}
      \label{fig:thm:cyclic_permutation_characterization}
    \end{figure}

    \thmitem{thm:cyclic_permutation_characterization/generating} If \( s_1, \ldots, s_m \) are distinct numbers between \( 1 \) and \( n \), then \eqref{eq:thm:def:cyclic_permutation/cycle} defines a cycle of length \( m \).
  \end{thmenum}
\end{proposition}
\begin{comments}
  \item Some authors like \textcite[15]{Knapp2016BasicAlgebra} and \textcite[sec. 4.3]{Тыртышников2007ЛинейнаяАлгебра} define cycles as permutations satisfying \eqref{eq:thm:def:cyclic_permutation/cycle}.
\end{comments}
\begin{proof}
  \SubProofOf{thm:cyclic_permutation_characterization/cycle} If \( \sigma \) is the identity, this is trivial. Suppose that \( \sigma \) is a cycle of length \( m > 1 \). We will prove the individual cases in \eqref{eq:thm:def:cyclic_permutation/cycle}:
  \begin{itemize}
    \item \Cref{thm:pigeonhole_principle/simple} implies that the value \( \sigma(s_m) \) must be among \( s_1, \ldots, s_m \).

    If \( \sigma(s_m) = s_i \) for \( i > 1 \), then \( \sigma(s_m) = s_i = \sigma(s_{i-1}) \), contradicting the injectivity of \( \sigma \).

    Hence, it follows that \( \sigma(s_m) = s_1 \).

    \item For any \( i < m \), by definition \( \sigma(s_i) = s_{i + 1} \).

    \item If \( q \) is not in the nontrivial orbit of \( \sigma \), it is the only element of its own orbit, and hence \( \sigma(q) = q \).
  \end{itemize}

  \SubProofOf{thm:cyclic_permutation_characterization/generating} Trivial.
\end{proof}

\begin{remark}\label{rem:cycle_notation}
  Having \cref{thm:cyclic_permutation_characterization} in mind, we can introduce a simplified notation for \hyperref[def:cyclic_permutation]{cyclic permutations}.

  Rather than the elaborate definition \eqref{eq:thm:def:cyclic_permutation/cycle}, we can write the corresponding cycle as
  \begin{equation*}
    \sigma = \cycle{ s_1, \ldots, s_m }.
  \end{equation*}

  We prefer denoting the identity by \( \id \), even though, for any number \( s \), the notation \( \cycle{ s } \) is just as valid.
\end{remark}

\begin{example}\label{ex:def:cyclic_permutation}
  We list examples of \hyperref[def:cyclic_permutation]{cyclic permutations}:
  \begin{thmenum}
    \thmitem{ex:def:cyclic_permutation/id} As already mentioned, the identity is the unique degenerate cycle of length \( 1 \), and can be expressed via the cycle notation \( \cycle{ n } \) for any number \( n \).

    \thmitem{ex:def:cyclic_permutation/123} Consider the permutation \( \cycle{ 1, 2, 3 } \) from \( S_3 \). Writing it via the permutation notation from \cref{def:symmetric_group}, we obtain
    \begin{equation*}
      \begin{pmatrix}
        1 & 2 & 3 \\
        2 & 3 & 1
      \end{pmatrix}.
    \end{equation*}

    Furthermore,
    \begin{equation*}
      \cycle{ 1, 2, 3 } = \cycle{ 1, 3 } \bincirc \cycle{ 1, 2 },
    \end{equation*}
    which generalizes to \cref{thm:cycle_transposition_decomposition}.

    We will discuss this group in detail in \cref{ex:s3}.

    \thmitem{ex:def:cyclic_permutation/diamond} Consider the permutation \eqref{eq:ex:def:symmetric_group/diamond} from \cref{ex:def:symmetric_group/diamond}. It has two orbits --- \( \set{ 1, 3, 5 } \) and \( \set{ 2, 4, 6 } \). The two orbits correspond to the cycles \( \cycle{ 1, 3, 5 } \) and \( \cycle{ 2, 4, 6 } \). Composing these cycles in any order gives \eqref{eq:ex:def:symmetric_group/diamond}.

    This is generalized in \cref{thm:permutation_decomposition_into_disjoint_cycles}.

    Furthermore, from \cref{fig:ex:def:symmetric_group/diamond} it becomes clear that
    \begin{equation*}
      \cycle{ 1, 3, 5 } = \cycle{ 3, 5, 1 } = \cycle{ 5, 1, 3 }.
    \end{equation*}

    This is generalized in \cref{thm:cyclic_permutation_cyclic_shift}.
  \end{thmenum}
\end{example}

\begin{definition}\label{def:shifted_remainder}\mimprovised
  Under our integer remainder function from \fullref{alg:integer_division}, \( 0 \leq \rem(n, m) < m \). Since we often start indexing from \( 1 \) rather than \( 0 \), it will be convenient to use the \enquote{shifted remainder}
  \begin{equation}\label{eq:def:shifted_remainder}
    \rem_1(n, m) \coloneqq \rem(n - 1, m) + 1,
  \end{equation}
  where \( 1 \leq \rem_1(n, m) \leq m \).
\end{definition}

\begin{definition}\label{def:cyclic_shift}\mimprovised
  We define the \term{right cyclic shift} of a list
  \begin{equation*}
    a_1, a_2, \ldots, a_n
  \end{equation*}
  by an integer \( c \) as the list
  \begin{equation*}
    a_{k_1}, \ldots, a_{k_n},
  \end{equation*}
  where \( k_i \coloneqq \rem_1(i + c, n) \) for \( i = 1, \ldots, n \).

  We denote by \( S^c \) the operator that sends a list of length \( n \) to its right shift by \( c \).

  We define left cyclic shifts as right cyclic shifts by \( -c \).
\end{definition}
\begin{comments}
  \item Cyclic shifts adapt to finite lists the shift operators from infinite and doubly-infinite sequences, which we define in \cref{def:sequence_shift_operator}, as well as digit shifts, which we define in \cref{def:digit_shift}.
\end{comments}

\begin{proposition}\label{thm:cyclic_shift_is_action}
  Fix a subset \( L \) of \( n \)-tuples that is closed under \hyperref[def:cyclic_shift]{cyclic shifts}.

  Then the \hyperref[def:cyclic_group]{cyclic group} \( C_n \) \hyperref[def:group_action]{acts} on \( S \) by sending \( a^c \) to a right shift by \( c \).
\end{proposition}
\begin{comments}
  \item Since \( C_n \) is commutative, this is simultaneously a left and a right action.
\end{comments}
\begin{proof}
  We must show that \( S^d \bincirc S^c = S^{d + c} \).

  For a concrete tuple \( (x_1, \ldots, x_n) \), we have
  \begin{equation*}
    S^d(S^c(x_1, \ldots, x_n)) = S^d(x_{k_1}, \ldots, x_{k_n}) = (x_{m_1}, \ldots, x_{m_n}),
  \end{equation*}
  where \( k_i \coloneqq \rem_1(c + i, n) = \rem(c + i - 1, n) + 1 \) and \( m_i \coloneqq \rem_1(d + k_i, n) = \rem(d + k_i - 1, n) + 1 \) for every index \( i \).

  We will discuss modular arithmetic in more detail in \fullref{sec:abelian_groups}; until then, we must prove from first principles whatever properties we need.

  Fix an index \( i \). By definition of remainder, there exist integers \( e \) such that
  \begin{equation*}
    c + i - 1 = ne + (k_i - 1)
  \end{equation*}
  and \( f \) such that
  \begin{equation*}
    d + k_i - 1 = nf + (m_i - 1).
  \end{equation*}

  We have \( k_i = c + i - ne \), thus
  \begin{equation*}
    d + (c + i - ne) - 1 = nf + (m_i - 1),
  \end{equation*}
  which we can transform to
  \begin{equation*}
    (d + c) + i - 1 = n(e + f) + (m_i - 1).
  \end{equation*}

  Therefore, \( m_i = \rem((d + c) + i - 1, n) + 1 = \rem_1((d + c) + i, n) \).

  We have thus demonstrated that \( S^d \bincirc S^c \) is a cyclic shift by \( d + c \).
\end{proof}

\begin{proposition}\label{thm:cyclic_permutation_cyclic_shift}
  \hyperref[def:cyclic_permutation]{Cyclic permutations} are invariant under \hyperref[def:cyclic_shift]{cyclic shifts} of the elements in their notation.

  More concretely, consider a cycle \( \cycle{ s_1, \cdots, s_m } \) from \( S_n \) and let \( k_1, \ldots, k_m \) be the coefficients of some cyclic shift of \( s_1, \ldots, s_n \). Then
  \begin{equation*}
    \cycle{ s_1, \cdots, s_m } = \cycle{ s_{k_1}, \cdots, s_{k_m} }.
  \end{equation*}
\end{proposition}
\begin{comments}
  \item If we visualize the cycle as in \cref{fig:thm:cyclic_permutation_characterization}, cyclic shifting corresponds to \hyperref[def:rigid_motion/rotation]{plane rotation}.
\end{comments}
\begin{proof}
  Obvious from the definition of cycle.
\end{proof}

\begin{definition}\label{def:disjoint_cycle}\mcite[214]{Aluffi2009Algebra}
  We say that two \hyperref[def:cyclic_permutation]{cycles} on \( S_n \) are \term[ru=независимые (\cite[sec. 4.3]{Тыртышников2007ЛинейнаяАлгебра})]{disjoint} if their nontrivial \hyperref[def:group_action_orbit]{orbits} are disjoint.
\end{definition}
\begin{comments}
  \item More simply, the cycles \( \cycle{ s_1, \ldots, s_m } \) and \( \cycle{ r_1, \ldots, r_l } \) are disjoint if the sets \( \set{ s_1, \ldots, s_m } \) and \( \set{ r_1, \ldots, r_l } \) are disjoint.
\end{comments}

\begin{proposition}\label{thm:disjoint_cycles_commute}
  \hyperref[def:disjoint_cycle]{Disjoint cycles} commute under composition.
\end{proposition}
\begin{proof}
  Trivial.
\end{proof}

\begin{proposition}\label{thm:permutation_decomposition_into_disjoint_cycles}
  Let \( \sigma \) be an arbitrary permutation in \( S_n \) and let \( O_1, \ldots, O_p \) be the nontrivial orbits of \( \sigma \).

  To each nontrivial orbit \( O_i \) there corresponds some cycle \( \cycle{ s_{i,1}, \ldots, s_{i,m_i} } \). Then
  \begin{equation*}
    \sigma = \cycle{ s_{1,1}, \ldots, s_{1,m_1} } \bincirc \cdots \bincirc \cycle{ s_{p,1}, \ldots, s_{p,m_p} }.
  \end{equation*}

  This gives us a decomposition of \( \sigma \) into pairwise \hyperref[def:disjoint_cycle]{disjoint} cycles. \Cref{thm:disjoint_cycles_commute} implies that we can rearrange them, forming other decompositions.
\end{proposition}
\begin{proof}
  Trivial.
\end{proof}

\begin{remark}\label{rem:identity_in_permutation_decomposition}
  Consider the permutation
  \begin{equation*}
    \begin{pmatrix}
      1 & 2 & 3 & 4 & 5 & 6 & 7 \\
      1 & 3 & 2 & 4 & 6 & 5 & 7
    \end{pmatrix}
  \end{equation*}
  which exchanges \( 2 \) with \( 3 \) and \( 5 \) with \( 6 \). The two nontrivial orbits are thus \( \set{ 2, 3 } \) and \( \set{ 5, 6 } \), and \cref{thm:permutation_decomposition_into_disjoint_cycles} implies that the permutation can be decomposed as
  \begin{equation*}
    \cycle{ 2, 3 } \bincirc \cycle{ 5, 6 }.
  \end{equation*}

  We can, following authors like \textcite[49]{Jacobson1985BasicAlgebraI} and \textcite[118]{Rotman2015AdvancedModernAlgebraPart1}, for completeness write the above as
  \begin{equation*}
     \cycle{ 1 } \bincirc \cycle{ 2, 3 } \bincirc \cycle{ 4 } \bincirc \cycle{ 5, 6 } \bincirc \cycle{ 7 }.
  \end{equation*}

  This may help visually distinguishing the above permutation from the analogous permutation in \( S_6 \), which simply does not act on \( 7 \).

  Nevertheless, to avoid clutter, we generally prefer not writing one-cycles in cycle decompositions.
\end{remark}

\paragraph{Transpositions}

\begin{definition}\label{def:transposition}\mcite[219]{Aluffi2009Algebra}
  A \term{transposition} is a \hyperref[def:cyclic_permutation]{cyclic permutation} of length \( 2 \).
\end{definition}

\begin{proposition}\label{thm:transpositions_are_involutions}
  Every \hyperref[def:transposition]{transposition} is an \hyperref[def:morphism_invertibility/involution]{involution}.
\end{proposition}
\begin{proof}
  Trivial.
\end{proof}

\begin{proposition}\label{thm:cycle_transposition_decomposition}
  Every nontrivial cycle \( \cycle{ s_1, \cdots, s_m } \) can be decomposed into the product of \hyperref[def:transposition]{transpositions}
  \begin{equation*}
    \cycle{ s_1, \cdots, s_m } = \cycle{ s_1, s_m } \bincirc \cycle{ s_1, s_{m-1} } \bincirc \cdots \bincirc \cycle{ s_1, s_2 }.
  \end{equation*}
\end{proposition}
\begin{proof}
  We will use induction by \( m \). The proposition is trivial for transpositions (\( m = 2 \)), hence suppose that it holds for \( m - 1 \) and consider the transposition \( \cycle{ s_1, \cdots, s_m } \).

  We have, by the inductive hypothesis,
  \begin{equation*}
    \cycle{ s_1, \cdots, s_{m - 1} }
    =
    \cycle{ s_1, s_{m - 1} } \bincirc \cycle{ s_1, s_{m-1} } \bincirc \cdots \bincirc \cycle{ s_1, s_2 }.
  \end{equation*}

  By definition of cycle, \( s_{m-1} \) and \( s_1 \) get swapped, therefore
  \begin{equation*}
    \cycle{ s_1, \cdots, s_m } = \cycle{ s_1, s_m } \bincirc \cycle{ s_1, \cdots, s_{m - 1} }.
  \end{equation*}

  This completes the induction.
\end{proof}

\begin{lemma}\label{thm:permutation_parity_correctness}
  If a \hyperref[def:symmetric_group]{permutation} \( \sigma \in S_n \) can be decomposed into \hyperref[def:transposition]{transpositions} as both
  \begin{equation}\label{eq:thm:permutation_parity_correctness/s}
    \sigma = \underbrace{\cycle{ s_1, s_2 } \bincirc \cycle{ s_3, s_4 } \bincirc \cdots \bincirc \cycle{ s_{2m-1}, s_{2m} }}_{m \T*{transpositions}}
  \end{equation}
  and
  \begin{equation}\label{eq:thm:permutation_parity_correctness/r}
    \sigma = \underbrace{\cycle{ r_1, r_2 } \bincirc \cycle{ r_3, r_4 } \bincirc \cdots \bincirc \cycle{ r_{2l-1}, r_{2l} }}_{l \T*{transpositions}},
  \end{equation}
  then \( m - l \) is an even number.
\end{lemma}
\begin{comments}
  \item The gist of the proposition is discussed in \cref{ex:thm:permutation_parity_correctness}.
\end{comments}
\begin{proof}
  We will use induction on \( m \). First consider the base case \( m = 0 \). Then \( \sigma \) is the identity. Hence, every transposition in \eqref{eq:thm:permutation_parity_correctness/r} should be present twice so that its action cancels out. Therefore, \( l \) is an even number.

  Now suppose that the statement holds for \( m - 1 \). Add (compose on the right) the last transposition \( \cycle{ s_{2m-1}, s_{2m} } \) of \eqref{eq:thm:permutation_parity_correctness/s} to both \eqref{eq:thm:permutation_parity_correctness/s} and \eqref{eq:thm:permutation_parity_correctness/r}. The obtained permutations are obviously equal. Furthermore, since \( \cycle{ s_{2m-1}, s_{2m} } \) is its own inverse, we can just as well remove \( \cycle{ s_{2m-1}, s_{2m} } \) from \eqref{eq:thm:permutation_parity_correctness/s} to obtain a decomposition of \( \sigma \bincirc \cycle{ s_{2m-1}, s_{2m} } \) into \( m - 1 \) (rather than \( m + 1 \)) transpositions.

  By the inductive hypothesis, \( (m - 1) - (l + 1) = m - l - 2 \) is an even number. Therefore, \( m - l \) is also an even number.
\end{proof}

\begin{example}\label{ex:thm:permutation_parity_correctness}
  We will demonstrate \cref{thm:permutation_parity_correctness}. Consider the permutation \( \sigma \) from \cref{ex:def:symmetric_group/diamond}. We have shown in \cref{ex:def:cyclic_permutation/diamond} that
  \begin{equation*}
    \sigma = \cycle{ 1, 3, 5 } \bincirc \cycle{ 2, 4, 6 }.
  \end{equation*}

  \Cref{thm:cycle_transposition_decomposition} implies that
  \begin{equation*}
    \sigma = \cycle{ 1, 5 } \bincirc \cycle{ 1, 3 } \bincirc \cycle{ 2, 4 } \bincirc \cycle{ 2, 6 }.
  \end{equation*}

  We can add any transposition to this decomposition, but we must also reverse its action to obtain \( \eqref{eq:ex:def:symmetric_group/diamond} \) again. For example, \( \sigma \bincirc \cycle{ 1, 2 } \) is not equal to \( \sigma \), but
  \begin{equation*}
    \sigma \bincirc \cycle{ 1, 2 } \bincirc \cycle{ 1, 2 } = \sigma.
  \end{equation*}

  Let us try to find another decomposition of \( \sigma \). First note that
  \begin{equation*}
    \sigma \bincirc \cycle{ 1, 2 }
    =
    \begin{pmatrix}
      1 & 2 & 3 & 4 & 5 & 6 \\
      3 & 4 & 5 & 6 & 2 & 1
    \end{pmatrix}
    =
    \cycle{ 1, 3, 5, 2, 4, 6 }
    \reloset {\ref{thm:cyclic_permutation_cyclic_shift}} =
    \cycle{ 3, 5, 2, 4, 6, 1 }.
  \end{equation*}

  We choose to start with \( 3 \) so that no two transpositions in our final representation are equal. \Cref{thm:cycle_transposition_decomposition} implies that
  \begin{equation*}
    \sigma \bincirc \cycle{ 1, 2 }
    =
    \underbrace{\cycle{ 3, 1 }}_{\cycle{ 1, 3 }} \bincirc \cycle{ 3, 6 } \bincirc \cycle{ 3, 4 } \bincirc \underbrace{\cycle{ 3, 2 }}_{\cycle{ 2, 3 }} \bincirc \cycle{ 3, 5 }.
  \end{equation*}

  Then
  \begin{equation*}
    \sigma
    \reloset {\ref{thm:transpositions_are_involutions}} =
    \sigma \bincirc \cycle{ 1, 2 } \bincirc \cycle{ 1, 2 }
    =
    \cycle{ 1, 3 } \bincirc \cycle{ 3, 6 } \bincirc \cycle{ 3, 4 } \bincirc \cycle{ 2, 3 } \bincirc \cycle{ 3, 5 } \bincirc \cycle{ 1, 2 }.
  \end{equation*}

  These transpositions are not disjoint, unlike those from our initial decomposition.
\end{example}

\begin{proposition}\label{thm:permutation_decomposition_existence}
  Every \hyperref[def:symmetric_group]{permutation} from \( S_n \) can be decomposed into a product of \hyperref[def:cyclic_permutation]{transpositions}.
\end{proposition}
\begin{comments}
  \item It follows from the discussion in \cref{ex:thm:permutation_parity_correctness} that this decomposition is not unique.
\end{comments}
\begin{proof}
  By \cref{thm:permutation_decomposition_into_disjoint_cycles}, the permutation can be decomposed into a (perhaps nullary) product of cycles. By \cref{thm:permutation_parity_correctness}, each of those cycles can be decomposed into a product of transpositions. Hence, at least one decomposition into transpositions exists for every permutation.
\end{proof}

\begin{definition}\label{def:permutation_parity}\mimprovised
  We say that a \hyperref[def:symmetric_group]{permutation} from \( S_n \) is \term{even} if it can be decomposed into an even number of \hyperref[def:cyclic_permutation]{transpositions}. Otherwise, we call the permutation \term{odd}.

  We correspondingly define the \term{sign} of a permutation as \( 1 \) for even permutations and as \( -1 \) for odd permutation.
\end{definition}
\begin{defproof}
  \Cref{thm:permutation_decomposition_existence} implies that at least one such decomposition exists.

  \Cref{thm:permutation_parity_correctness} implies that if one such decomposition has an even (resp. odd) number of transpositions, every other decomposition also has an even (resp. odd) number of transpositions.
\end{defproof}

\begin{proposition}\label{thm:permutation_product_parity}
  The composition of two \hyperref[def:symmetric_group]{permutations} from \( S_n \) is \hyperref[def:permutation_parity]{odd} if and only if they have different signs.
\end{proposition}
\begin{proof}
  When the permutations \( \sigma \) and \( \rho \) are decomposed into \( n \) and \( m \) transpositions, correspondingly, their product has \( n + m \) transpositions.

  Hence, their product is even or odd in accordance with \cref{thm:integer_sum_parity}.
\end{proof}

\begin{proposition}\label{thm:permutation_sign_homomorphism}
  The \hyperref[def:permutation_parity]{permutation sign} \( \sgn: S_n \to \BbbZ \) is a group homomorphism (with respect to multiplication in \( \BbbZ \)).
\end{proposition}
\begin{proof}
  We saw in \cref{thm:permutation_product_parity} that the composition \( \sigma \bincirc \pi \) is even if and only if \( \sigma \) and \( \pi \) have the same signs. Thus,
  \begin{equation*}
    \sgn(\sigma \bincirc \pi) = \sgn(\sigma) \cdot \sgn(\pi),
  \end{equation*}
  as desired.
\end{proof}

\begin{definition}\label{def:cayley_table}\mcite[150]{Rotman2015AdvancedModernAlgebraPart1}
  The \term[ru=таблица умножения (\cite[135]{Шафаревич1999ОсновныеПонятияАлгебры})]{multiplication table} or \term[ru=таблица Кэли (\cite[135]{Шафаревич1999ОсновныеПонятияАлгебры})]{Cayley table} for a finite \hyperref[def:semigroup]{(semi)group} \( G = \set{ a_1, \ldots, a_n } \), is an \( n \times n \) matrix whose \( (i, j) \)-th entry is \( a_i a_j \).

  More generally, we can define a Cayley table for every binary operation on a finite set, however calling it a multiplication table is not appropriate in general.
\end{definition}

\begin{example}\label{ex:s3}
  The \hyperref[def:symmetric_group]{symmetric group} \( S_3 \) contains the following \hyperref[def:symmetric_group]{permutations}:
  \begin{equation*}
    S_3
    \coloneqq
    \set[\vast]
    {
      \underbrace
        {
          \begin{pmatrix}
            1 & 2 & 3 \\
            1 & 2 & 3
          \end{pmatrix}
        }_{
          \id
        },
      \underbrace
        {
          \begin{pmatrix}
            1 & 2 & 3 \\
            2 & 1 & 3
          \end{pmatrix}
        }_{
          \cycle{ 1, 2 }
        },
      \underbrace
        {
          \begin{pmatrix}
            1 & 2 & 3 \\
            2 & 3 & 1
          \end{pmatrix}
        }_{
          \cycle{ 1, 2, 3 }
        },
      \underbrace
        {
          \begin{pmatrix}
            1 & 2 & 3 \\
            3 & 2 & 1
          \end{pmatrix}
        }_{
          \cycle{ 1, 3 }
        },
      \underbrace
        {
          \begin{pmatrix}
            1 & 2 & 3 \\
            3 & 1 & 2
          \end{pmatrix}
        }_{
          \cycle{ 1, 3, 2 }
        },
      \underbrace
        {
          \begin{pmatrix}
            1 & 2 & 3 \\
            1 & 3 & 2
          \end{pmatrix}
        }_{
          \cycle{ 2, 3 }
        }
    }
  \end{equation*}

  We will now construct a \hyperref[def:cayley_table]{Cayley table} to show that \( S_3 \) is not commutative. This will also help us build examples based on \( S_3 \) later on.

  The entire table can be filled using the following techniques:
  \begin{itemize}
    \item The identity permutation \( \id \) leaves the other multiplicand unchanged.
    \item The cycles \( \cycle{ 1, 2, 3 } \) and \( \cycle{ 1, 3, 2 } \) are inverses of each other, and \cref{thm:transpositions_are_involutions} implies that all others are \hyperref[def:morphism_invertibility/involution]{involutions}.
    \item \Cref{thm:cycle_transposition_decomposition} implies that \( \cycle{ 1, 2, 3 } = \cycle{ 1, 3 } \bincirc \cycle{ 1, 2 } \) and \( \cycle{ 1, 3, 2 } = \cycle{ 1, 2 } \bincirc \cycle{ 1, 3 } \).
    \item The last two points together imply
    \begin{equation*}
      \cycle{ 1, 2, 3 } \bincirc \cycle{ 1, 2 }
      =
      \cycle{ 1, 3 } \bincirc \cycle{ 1, 2 } \bincirc \cycle{ 1, 2 }
      =
      \cycle{ 1, 3 }.
    \end{equation*}

    \item We can also utilize cyclic shifts as justified by \cref{thm:cyclic_permutation_cyclic_shift}:
    \begin{equation*}
      \cycle{ 1, 2 } \bincirc \cycle{ 2, 3 }
      =
      \cycle{ 2, 1 } \bincirc \cycle{ 2, 3 }
      =
      \cycle{ 2, 3, 1 }
      =
      \cycle{ 1, 2, 3 }
    \end{equation*}

    Other cases can be handled similarly.
  \end{itemize}

  All in all, \cref{tab:ex:s3} provides all values of \( \tau \bincirc \sigma \).

  \begin{table}[!ht]
    \begin{equation*}
      \begin{array}{l !{\quad} *{6}{l}}
        \toprule
        \multicolumn{1}{c}{\tau \bincirc \sigma} & \multicolumn{6}{c}{\tau \bincirc \sigma} \\
        \cmidrule{2-7}
        \multicolumn{1}{c}{\tau}                 & \id               & \cycle{ 1, 2 }    & \cycle{ 1, 3 }    & \cycle{ 2, 3 }    & \cycle{ 1, 2, 3 } & \cycle{ 1, 3, 2 } \\
        \midrule
        \id                                      & \id               & \cycle{ 1, 2 }    & \cycle{ 1, 3 }    & \cycle{ 2, 3 }    & \cycle{ 1, 2, 3 } & \cycle{ 1, 3, 2 } \\
        \cycle{ 1, 2 }                           & \cycle{ 1, 2 }    & \id               & \cycle{ 1, 3, 2 } & \cycle{ 1, 2, 3 } & \cycle{ 2, 3 }    & \cycle{ 1, 3 }    \\
        \cycle{ 1, 3 }                           & \cycle{ 1, 3 }    & \cycle{ 1, 2, 3 } & \id               & \cycle{ 1, 3, 2 } & \cycle{ 1, 2 }    & \cycle{ 2, 3 }    \\
        \cycle{ 2, 3 }                           & \cycle{ 2, 3 }    & \cycle{ 1, 3, 2 } & \cycle{ 1, 2, 3 } & \id               & \cycle{ 1, 3 }    & \cycle{ 1, 2 }    \\
        \cycle{ 1, 2, 3 }                        & \cycle{ 1, 2, 3 } & \cycle{ 1, 3 }    & \cycle{ 2, 3 }    & \cycle{ 1, 2 }    & \cycle{ 1, 3, 2 } & \id               \\
        \cycle{ 1, 3, 2 }                        & \cycle{ 1, 3, 2 } & \cycle{ 2, 3 }    & \cycle{ 1, 2 }    & \cycle{ 1, 3 }    & \id               & \cycle{ 1, 2, 3 } \\
        \bottomrule
      \end{array}
    \end{equation*}
    \caption{\hyperref[def:cayley_table]{Multiplication table} for the \hyperref[def:symmetric_group]{symmetric group} \( S_3 \)}\label{tab:ex:s3}
  \end{table}
\end{example}

\begin{proposition}\label{thm:s3_and_d6}
  The \hyperref[def:symmetric_group]{symmetric group} \( S_3 \) is isomorphic to the \hyperref[def:dihedral_group]{dihedral group} \( D_6 \).
\end{proposition}
\begin{proof}
  Define an isomorphism as follows:
  \begin{equation*}
    \begin{aligned}
      &\varphi: S_3 \to D_6, \\
      &\varphi(\sigma) \coloneqq \begin{cases}
        a,   \sigma = \cycle{ 1, 2, 3 }, \\
        a^2, \sigma = \cycle{ 1, 3, 2 }, \\
        b,   \sigma = \cycle{ 1, 2 }, \\
      \end{cases}
    \end{aligned}
  \end{equation*}

  From the multiplication table in \cref{ex:s3} we can see that \( \varphi(a \cdot a) \) indeed coincides with \( \varphi(a^2) \).

  We must now verify the three equations from \eqref{eq:def:dihedral_group}.
  \begin{itemize}
    \item \( \varphi(a^3) = \varphi(a^2) \cdot \varphi(a) = \cycle{ 1, 3, 2 } \bincirc \cycle{ 1, 2, 3 } = \id \), as desired.
    \item \( \varphi(b^2) = \varphi(b)^2 = \cycle{ 1, 2 }^2 = \id \), as desired.
    \item \( \varphi(bab) = \varphi(ba) \cdot \varphi(b) = \cycle{ 2, 3 } \bincirc \cycle{ 1, 2, 3 } = \cycle{ 1, 3, 2 } = \varphi(a^2) \), as desired (since \( a^2 = a^{-1} \)).
  \end{itemize}
\end{proof}

\paragraph{Alternating groups}

\begin{definition}\label{def:alternating_group}\mcite[def. IV.4.12]{Aluffi2009Algebra}
  We define the \term[bg=алтернативна група (\cite[379]{Обрешков1962ВисшаАлгебра})]{alternating group} \( A_n \) on \( n \) letters as the subgroup of all \hyperref[def:permutation_parity]{even permutation} in the \hyperref[def:symmetric_group]{symmetric group} \( S_n \).
\end{definition}
\begin{comments}
  \item The alternating group \( A_n \) is \hyperref[def:normal_subgroup]{normal} in \( S_n \) as the kernel of the \hyperref[def:permutation_parity]{permutation sign} homomorphism.
\end{comments}

\begin{example}\label{ex:s3_and_a3}
  In the \hyperref[def:symmetric_group]{symmetric group} \( S_3 \) discussed in \cref{ex:s3}, the \hyperref[def:alternating_group]{alternating group} \( A_3 \) consists of:
  \begin{itemize}
    \item The identity, which is a product of zero transpositions.
    \item The even-length cycle \( \cycle{ 1, 2, 3 } = \cycle{ 1, 3 } \bincirc \cycle{ 1, 2 } \).
    \item The even-length cycle \( \cycle{ 1, 3, 2 } = \cycle{ 1, 2 } \bincirc \cycle{ 1, 3 } \).
  \end{itemize}
\end{example}

\begin{proposition}\label{thm:group_epimorphisms_are_surjective}\mcite[exerc. I.5.5]{MacLane1998CategoryTheory}
  Every \hyperref[def:fol_epimorphism]{epimorphism} in \hyperref[def:group/category]{\( \cat{Grp} \)} is \hyperref[def:function_invertibility/surjective]{surjective}.
\end{proposition}
\begin{proof}
  Let \( \varphi: G \to H \) be an epimorphism and suppose that it is not surjective. Let \( M \) be the \hyperref[def:normal_closure]{normal closure} \( \img \varphi \).

  Aiming at a contradiction, suppose that \( M \) has \hyperref[def:subgroup_index]{index} \( 2 \) in \( H \), consider the quotient map \( \pi: H \to H / M \) and the constant map \( c(h) \coloneqq M \). Then
  \begin{equation*}
    \pi \bincirc \varphi = c \bincirc \varphi.
  \end{equation*}

  Since \( \varphi \) is an epimorphism, we have \( \pi = c \). But we have deliberately taken \( \pi \) and \( c \) so that \( \pi \neq c \). The obtained contradiction shows that \( M \) must have an index greater than \( 2 \).

  Thus, the index of \( M \) in \( G \) is greater than \( 2 \). Let \( M \), \( uM \) and \( vM \) be different cosets. Define \( \sigma: H \to H \) as the \hyperref[def:symmetric_group]{permutation} on \( H \) that exchanges \( xu \) with \( xv \) for every \( x \in M \). \Cref{thm:group_conjugation_action} implies that the following is a homomorphism:
  \begin{equation*}
    \begin{aligned}
      &\theta: H \to S_H \\
      &\theta(h) \coloneqq \sigma^{-1} \bincirc \psi(h) \bincirc \sigma.
    \end{aligned}
  \end{equation*}

  Consider also another homomorphism,
  \begin{equation*}
    \begin{aligned}
      &\psi: H \to S_H \\
      &\psi(h) \coloneqq (x \mapsto hx),
    \end{aligned}
  \end{equation*}
  where \( S_H \) is the \hyperref[def:symmetric_group]{symmetric group}. This is indeed a homomorphism by \fullref{thm:cayleys_theorem}.

  Since \( \sigma \) fixes the members of \( M \) in-place, we have \( \theta(h)\restr_M = \psi(h)\restr_M \). Since \( M \) contains the image of \( \varphi \), this implies
  \begin{equation*}
    \psi \bincirc \varphi = \theta \bincirc \varphi.
  \end{equation*}

  Since \( \varphi \) is an epimorphism, we have \( \psi = \theta \). But we have deliberately constructed \( \psi \) and \( \theta \) such that \( \psi \neq \theta \). The obtained contradiction shows that \( \img \varphi \) cannot be a strict subgroup of \( G \). Therefore, \( \varphi \) must be surjective.
\end{proof}

\begin{proposition}\label{thm:sn_semidirect_product}
  The \hyperref[def:symmetric_group]{symmetric group} \( S_n \) decomposes into the \hyperref[def:internal_semidirect_product]{semidirect product}
  \begin{equation*}
    S_n = A_n \rtimes \braket{ \cycle{ 1, 2 } }.
  \end{equation*}
\end{proposition}
\begin{proof}
  Follows from \cref{thm:index_2_semidirect_product}.
\end{proof}

\begin{proposition}\label{thm:an_commutativity}
  The \hyperref[def:alternating_group]{alternating group} \( A_n \) is \hyperref[def:abelian_group]{abelian} if and only if \( n \leq 3 \).
\end{proposition}
\begin{proof}
  The groups \( A_1 \) and \( A_2 \) are trivial, and commutativity for the even-length cycles in \( S_3 \) is shown in \cref{ex:s3}. Therefore, \( A_n \) is abelian when \( n \leq 3 \).

  On the other hand, when \( n > 3 \) the cycles \( \cycle{ 1, 2, 3 } \) and \( \cycle{ 2, 3, 4 } \) do not commute because we have
  \begin{equation*}
    \underbrace
      {
        \begin{pmatrix}
          1 & 2 & 3 & 4 \\
          1 & 3 & 4 & 2
        \end{pmatrix}
      }_{\cycle{ 2, 3, 4 }}
    \bincirc
    \underbrace
      {
        \begin{pmatrix}
          1 & 2 & 3 & 4 \\
          2 & 3 & 1 & 4
        \end{pmatrix}
      }_{\cycle{ 1, 2, 3 }}
    =
    \underbrace
      {
        \begin{pmatrix}
          1 & 2 & 3 & 4 \\
          3 & 4 & 1 & 2
        \end{pmatrix}
      }_{\cycle{ 2, 4 } \bincirc \cycle{ 1, 3 }}
  \end{equation*}
  and
  \begin{equation*}
    \underbrace
      {
        \begin{pmatrix}
          1 & 2 & 3 & 4 \\
          2 & 3 & 1 & 4
        \end{pmatrix}
      }_{\cycle{ 1, 2, 3 }}
    \bincirc
    \underbrace
      {
        \begin{pmatrix}
          1 & 2 & 3 & 4 \\
          1 & 3 & 4 & 2
        \end{pmatrix}
      }_{\cycle{ 2, 3, 4 }}
    =
    \underbrace
      {
        \begin{pmatrix}
          1 & 2 & 3 & 4 \\
          2 & 1 & 4 & 3
        \end{pmatrix}
      }_{\cycle{ 3, 4 } \bincirc \cycle{ 1, 2 }}
  \end{equation*}
\end{proof}
